\section{The quartic domain and finite-order centres}\label{inf:sec-centres}

In this section $r$ ranges over the fixed compact interval
$[1/4,3/8]$. All finite-order constants are uniform on this interval
and for $|\theta|\le20$, including $\theta=0$ when analytic
extensions are discussed. The accuracy order is always fixed before
$p$ tends to infinity.

\subsection{The quartic profile}

The monic degree-two angular polynomial is
\begin{equation}\label{inf:H2}
 \mathsf H_2(x)=x^2-\frac{2(\ab+1)}{p+2}x
                 +\frac{\ab(\ab+1)}{(p+1)(p+2)}.
\end{equation}
Its homogeneous spherical degree is four. Define its self-coupling by
\begin{equation}\label{inf:seed-coupling}
 g_2=[\mathsf H_2]\mathsf H_2^2
 =\frac{\int \mathsf H_2^3\,d\mu_r}{\int \mathsf H_2^2\,d\mu_r}.
\end{equation}
The bracket denotes the coefficient in the monic orthogonal expansion.

\begin{lemma}[Quartic normalization]\label{inf:quartic-normalization}
Put $\varsigma=(r-1/2)^2\in[1/64,1/16]$. Then
\begin{equation}\label{inf:g2-exact}
 g_2=\frac{p-1}{(p+1)(p+5)}
 -\frac{4p^2(p^2-9p-16)\varsigma}{(p+1)(p+2)^2(p+4)(p+5)}.
\end{equation}
For $p\ge64$,
\begin{equation}\label{inf:g2-bounds}
 \frac9{16}<pg_2<1,\qquad
 \nrm{\mathsf H_2}_\infty<\frac9{16},\qquad
 \nrm{T_\varsigma}_\infty<1,\quad T_\varsigma(x)=\frac{\mathsf H_2(x)}{pg_2}.
\end{equation}
Also $g_2=4r(1-r)p^{-1}+O(p^{-2})$ uniformly in $r$.
\end{lemma}
\begin{proof}
The beta moments are $\int x^h\,d\mu_r=(\ab)_h/(p)_h$.
They give the two orthogonality conditions for \eqref{inf:H2}.
Multiplying the three-term recurrence twice gives
\[
 \int \mathsf H_2P_j^2\,d\mu_r
 =\mathfrak a_j^2+\mathfrak a_{j+1}^2+\mathfrak b_j^2
 -\frac{2(\ab+1)}{p+2}\mathfrak b_j
 +\frac{\ab(\ab+1)}{(p+1)(p+2)}.
\]
At $j=2$, substitute \eqref{inf:recurrence-a}--\eqref{inf:recurrence-b}
and $\ab=p/2+p(r-1/2)$. The terms linear in $r-1/2$ cancel;
the constant term is $(p-1)/((p+1)(p+5))$, and the coefficient of
$(r-1/2)^2$ is the second coefficient in \eqref{inf:g2-exact}.
Since $P_2$ is a multiple of $\mathsf H_2$, this gives $g_2$.

Write $g_2=c_2^{\mathrm{diag}}+d_2^{\mathrm{diag}}\varsigma$
with the coefficients in \eqref{inf:g2-exact}. For $p\ge64$,
$d_2^{\mathrm{diag}}<0$, $|pd_2^{\mathrm{diag}}|<4$, and
$pc_2^{\mathrm{diag}}\ge4032/4485>13/16$.
The last inequality persists for larger $p$ because the derivative of
$p(p-1)/((p+1)(p+5))$ has numerator $7p^2+10p-5>0$.
Thus $pg_2>13/16-4/16=9/16$, while $pg_2<pc_2^{\mathrm{diag}}<1$.
The two positive endpoint values of $\mathsf H_2$ are less than $9/16$,
using $\ab,\bb\le3p/4$. Its negative minimum has magnitude
$(\ab+1)(\bb+1)/((p+1)(p+2)^2)\le1/(4(p+1))$.
This proves \eqref{inf:g2-bounds}. Expansion of the rational formula
gives $1-4\varsigma=4r(1-r)$ as its leading coefficient.
\end{proof}

The physical quartic radius and its normalized version are
\begin{equation}\label{inf:quartic-domain}
 R+h_{\mathrm q}(x),\qquad
 h_{\mathrm q}=-\frac{2\chi}{g_2}\mathsf H_2,\qquad
 1+f_{\mathrm q}(x)=1-\frac{\theta}{2R^2}T_\varsigma(x).
\end{equation}
We denote the corresponding physical domain by $\Omega_{\mathrm q}$.
For real block parameters the same formula denotes its weighted
meridian domain. The scalar $\theta$ and the radius $R$ stay fixed
when $\varsigma$ varies.

\subsection{An analytic finite-order construction}

\begin{theorem}[Finite-order centres]\label{inf:centres}
Fix an integer $N\ge4$. For every prescribed finite radius
$R_{\mathrm{an}}>0$, put $U=\{x\in\C:|x|\le R_{\mathrm{an}}\}$.
There exist $C_{N,U},P_{N,U}<\infty$ such that the
following assertion holds. Suppose $p\in\tfrac12\Z$, $p\ge P_{N,U}$,
$J_p(R)=0$, $R\in[2p-1,2p+4]$, and
$0<|\theta|\le20$ with \eqref{inf:detuning-definition}.
For every $r\in[1/4,3/8]$ one can choose a polynomial $h_N(x)$ and
a finite sum $u_N$ of regular Helmholtz modes. For real block
parameters the equation means the weighted operator
\eqref{inf:invariant-laplacian}; at integer blocks it is the
Euclidean equation. These choices satisfy
$(\Delta+1)u_N=0$ exactly and
\begin{equation}\label{inf:centre-complex-defect}
 \sup_{x\in U}\left(|u_N(R+h_N(x),x)-1|
       +|\partial_\varrho u_N(R+h_N(x),x)|\right)
 \le C_{N,U}\theta^2p^{-N}.
\end{equation}
The number and degrees of the modes depend on $N$, not on $p,r,\theta$.
For every fixed polynomial coefficient norm and every fixed derivative
order $k$ on $U$, with constants $C_{N,U,k}$ and threshold
$P_{N,U,k}$,
\begin{align}
 h_N&=-\frac{2\chi}{g_2}\mathsf H_2+O_{N,U}(|\theta|p^{-3}),
 \label{inf:centre-profile}\\
 [\mathsf H_2]h_N&=-\frac{\theta}{16r(1-r)p}
                                  +O_{N,U}(|\theta|p^{-2}).
 \label{inf:centre-witness}
\end{align}
The radius $R_{\mathrm{an}}$ can be chosen arbitrarily large before
the constants and threshold are fixed. In particular the statement
applies as follows: fix $\rho'>2$ and $d>\rho'$, and then
choose one $R_{\mathrm{an}}>1+d$.
\end{theorem}

\begin{proof}
We construct a finite jet with more terms than the requested accuracy.
Put $\eps=1/p$, write $R=\rho/\eps$, and set
$Z=\rho^2$ and $S_0=4(1+3\eps+2\eps^2)$.
Equation~\eqref{inf:trace-ratio-exact} is equivalent to
\begin{equation}\label{inf:rho-equation}
 2Z(Z-S_0)+\theta\eps^2(1+\eps)(2Z-S_0)=0.
\end{equation}
At $(\eps,Z)=(0,4)$ its derivative in $Z$ is eight.
The analytic implicit-function theorem gives a unique root near four,
uniformly for $|\theta|\le21$. The supplied radius satisfies
$(R/p)^2\to4$ and, by \eqref{inf:trace-ratio-exact},
equation~\eqref{inf:rho-equation} with $Z=(R/p)^2$.
Thus $Z$ is this root for sufficiently large $p$. The other quadratic root is
$O(\theta\eps^2)$. Take the positive square root near two.
Then
\begin{equation}\label{inf:chi-analytic}
 \rho=2+3\eps+O(\eps^2),\qquad
 \chi=\frac{\theta\eps^2}{4\rho}
      =\frac18\theta\eps^2+O(\theta\eps^3).
\end{equation}
Here and below analyticity means on an open complex neighbourhood of
the stated compact real parameter set.

For finitely many angular indices, the boundary traces of regular
Bessel modes are computed from
\begin{equation}\label{inf:finite-q-recurrence}
 q_{-1}=1,\quad q_0=0,\quad
 q_{k+1}=\frac{2(1+k\eps)}\rho q_k-q_{k-1}.
\end{equation}
This finite recurrence is analytic in $\eps,\theta$.
Its limiting zero indices give the resonant angular set
$\mathcal I_{\mathrm{res}}=\{2,5,8,\ldots\}$.
For $j\notin\mathcal I_{\mathrm{res}}$ normalize the mode $f_j$ by $f_j(R)=1$;
for $j\in\mathcal I_{\mathrm{res}}$ normalize it by $f_j'(R)=1$.
The latter convention gives $f_2(R)=\chi$. For $j\ge5$ in
$\mathcal I_{\mathrm{res}}$,
\begin{equation}\label{inf:resonant-trace}
 f_j(R)=c_j\eps+O_j(\eps^2),\qquad
 c_j=\frac{(2j-1)(2j-4)}3>0.
\end{equation}
Indeed the first correction of \eqref{inf:finite-q-recurrence} at
$k=3h$ is $(-1)^h3h(h-1)\eps$, while the corresponding normal
trace tends to $(-1)^h$. Put $3h=2j-1$ to obtain
\eqref{inf:resonant-trace}. The change in $\rho$ due to $\theta$
starts at order $\eps^2$, so this coefficient is independent of
$\theta$. The nonresonant normalization denominators tend to $1$ or
$-1$ and are units in the analytic coefficient ring.

All collar derivatives are supplied by
\begin{equation}\label{inf:collar-ode}
 f_j''+\frac{(2-\eps)}{\rho+\eps y}f_j'
 +\left(1-\frac{2j(2j\eps^2+2\eps-2\eps^2)}
                    {(\rho+\eps y)^2}\right)f_j=0,
 \qquad y=\varrho-R.
\end{equation}
Its coefficients are analytic near $\eps=0$ on a fixed collar.
The ball mode has traces $u_0(R)=1$, $u'_0(R)=0$, $u''_0(R)=-1$.
The collar construction is a local analytic family.
At the actual inputs $J_p(R)=0$ its traces
are precisely those of the regular Bessel modes; uniqueness of the
ordinary differential equation identifies them on the collar and
then throughout their regular continuation.

We next record the angular filtration that makes the recursion finite.
For fixed $i,j$ write
\begin{equation}\label{inf:monic-product-valuation}
 \mathsf H_i\mathsf H_j=\sum_{k\le i+j}\mathfrak c_{ij}^k\mathsf H_k,\qquad
 \mathfrak c_{ij}^k=O\left(\eps^{\lceil(i+j-k)/2\rceil}\right).
\end{equation}
To prove this, use the monic recurrence
$\mathsf H_{j+1}=(x-\mathfrak b_j)\mathsf H_j-\mathfrak a_j^2\mathsf H_{j-1}$.
For fixed indices \eqref{inf:recurrence-a}--\eqref{inf:recurrence-b}
give $\mathfrak b_j-r=O(\eps)$ and
$\mathfrak a_j^2=O(\eps)$.
Induction shows that the coefficient of $(x-r)^k$ in $\mathsf H_j$ has
valuation at least $\lceil(j-k)/2\rceil$. Descending elimination
proves the same valuation for the inverse monic triangular
connection. Multiplying and composing these two connections proves
\eqref{inf:monic-product-valuation}, since a sum of ceilings is at
least the ceiling of the sum. In particular,
\begin{equation}\label{inf:seed-square}
 \mathsf H_2^2=\mathsf H_4+O(\eps)\mathsf H_3+g_2\mathsf H_2+O(\eps^2)\mathsf H_1+O(\eps^2)\mathsf H_0,
 \qquad g_2=4r(1-r)\eps+O(\eps^2).
\end{equation}

Write the formal field and boundary as
\[
 u=u_0+\sum_j a_jf_j\mathsf H_j,\qquad \varrho=R+h(x).
\]
Every coefficient of $\eps$ is a finite polynomial in $x$.
Eliminate $h$ from the radial Neumann condition first. Its derivative
in $h$ at zero is $u_0''(R)=-1$, so at each order the new coefficient
of $h$ is determined by the already known coefficients and the new
field coefficients, using the unit divisor $u_0''(R)=-1$.
In the Dirichlet equation eliminate $a_j$ for $j\notin\mathcal I_{\mathrm{res}}$.
Its diagonal linear coefficient is $f_j(R)=1$. Thus its new
coefficients are also determined by unit divisions.
This elimination includes the $j=0$ field coefficient.

Parametrize the remaining coefficients by
\begin{equation}\label{inf:resonant-scaling}
 a_2=\theta\eps \SeedA,\qquad
 a_j=\theta^2\eps^3 Z_j\quad(j\in\mathcal I_{\mathrm{res}},\ j\ge5).
\end{equation}
After the two eliminations the residual Dirichlet coefficients
$F_j$ have the form
\begin{align}
 \frac{F_2}{\theta^2\eps^3}
 &=\frac{\SeedA}{8}+2r(1-r)\SeedA^2+O(\eps),\label{inf:seed-equation}\\
 \frac{F_j}{\theta^2\eps^4}
 &=c_jZ_j+B_j(\SeedA,\theta,r)+O(\eps),
 \qquad j\in\mathcal I_{\mathrm{res}},\ j\ge5.\label{inf:other-equations}
\end{align}
Here $B_j$ is a finite polynomial and is zero except possibly at
$j=5,8$. These divisions are identities of formal series, including
at $\theta=0$.

The linear seed term is
$a_2\chi=\theta^2\eps^3 \SeedA/8+O(\theta^2\eps^4)$.
Eliminating $h$ gives quadratic Dirichlet term
$a_2^2\mathsf H_2^2/2$; its seed coefficient is the second term of
\eqref{inf:seed-equation} by \eqref{inf:seed-square}.
A pure cubic term has angular degree at most six; reaching index two
loses four degrees and contributes two further powers of $\eps$.
The leading nonresonant coefficient is $a_4=-a_2^2/2+O(\eps^3)$,
and its product with the seed has the same loss when projected to
index two. Thus neither contributes to the leading seed equation.
Quadratic terms have no nonseed resonant index. Cubic terms first
reach index five after losing one angular degree, which supplies one
extra $\eps$. At total order four the only possible new resonant
indices are five and eight. A new $Z_j$ contributes linearly through
\eqref{inf:resonant-trace}, giving $c_jZ_j$.
Its product with the seed has top index $j+2$, which is nonresonant;
reaching a resonant index at most $j$ loses at least two degrees and
therefore at least one extra $\eps$. Contributions through a newly
eliminated nonresonant coefficient have the same extra factor after
its unit division. At order zero in \eqref{inf:other-equations},
$Z_j$ enters linearly as $c_jZ_j$. Every other term has at least two
amplitudes or the factor $a_2\chi$, so it is divisible by $\theta^2$.

Choose the nonzero root
$\SeedA_0=-1/(16r(1-r))$ of \eqref{inf:seed-equation}.
The derivative of its left side with respect to $\SeedA$ at this root is
$-1/8$. After choosing $\SeedA_0$, determine
$Z_{j,0}=-B_j(\SeedA_0,\theta,r)/c_j$; these have finite support.
Work in the ring of formal
$\eps$ series whose coefficients are polynomials in $x$, with
coefficients analytic in $r,\theta$ on the indicated compact sets.
A series is required to have a finite polynomial coefficient at
every order.
For a fixed list of resonant amplitude coefficients the two unit
eliminations have linearization, in the variables $(h,a_{\rm nr})$,
\[
 (\dot h,\dot a_{\rm nr})\longmapsto
 \left(-\dot h+\sum_{j\notin\mathcal I_{\mathrm{res}}}
       f'_j(R)|_{\eps=0}\,\dot a_j\mathsf H_j|_{\eps=0},
       \dot a_{\rm nr}\right).
\]
Its inverse is triangular and sends a finite polynomial forcing
to finite polynomial coefficients. At each order, the remaining
terms involve only coefficients at lower orders. This proves
existence and uniqueness of those unit eliminations by ordinary
coefficient induction. It also proves preservation of finite
support: multiplication and the monic connection use only finite
polynomials, and the inverse displayed here adds no angular
indices not present in its forcing. Dependence of the exact
collar traces on $\eps$ contributes only previously determined
lower-order coefficients in this calculation.

After these eliminations let
$\mathcal F_2=F_2/(\theta^2\eps^3)$ and
$\mathcal F_j=F_j/(\theta^2\eps^4)$ for nonseed resonant $j$.
Their constant terms are the right sides of
\eqref{inf:seed-equation}--\eqref{inf:other-equations} without the
remainders. At $(\SeedA,Z)=(\SeedA_0,Z_0)$ their leading Jacobian is
\begin{equation}\label{inf:formal-Jacobian}
 \mathcal J_0(\dot{\SeedA},\dot Z)=
 \left(-\frac18\dot{\SeedA},
       \{\partial_{\SeedA}B_j(\SeedA_0,\theta,r)\dot{\SeedA}+c_j\dot Z_j\}_{j\ge5,
                                                    j\in\mathcal I_{\mathrm{res}}}\right).
\end{equation}
The off-diagonal column in this matrix has finite support,
since $B_j=0$ outside $j=5,8$. Each $c_j$ is positive and fixed.
Thus its inverse preserves finite support, using only the fixed
nonzero divisors $-1/8$ and $c_j$.

Suppose all coefficients of $\SeedA,Z$ below order $v\ge1$ have
already been selected and the reduced equations vanish below
that order. Substitute those truncated series with the unknown
order-$v$ coefficients set to zero, and denote the coefficients
of $\eps^v$ in $\mathcal F_2,\mathcal F_j$ by
$e_{2,v},e_{j,v}$. They are explicitly defined by finite
coefficient extraction from the collar Taylor formula, the
polynomial product connection, and the two unit eliminations.
Only finitely many $e_{j,v}$ are nonzero.
The new coefficients occur only through
\eqref{inf:formal-Jacobian}: every other occurrence is multiplied
by an additional $\eps$, or by a coefficient of positive order
of $\SeedA-\SeedA_0$ or $Z-Z_0$. Therefore the unique update is
\begin{equation}\label{inf:formal-update}
 \SeedA_v=8e_{2,v},\qquad
 Z_{j,v}=-\frac{e_{j,v}
       +\partial_{\SeedA}B_j(\SeedA_0,\theta,r)\SeedA_v}{c_j}.
\end{equation}
Set the other coefficients to zero. The update is finite, cancels
exactly the order-$v$ coefficient, and leaves every lower order
unchanged. Recompute the unit-eliminated coefficients to that
order. This completes an induction producing every finite jet.

All operations in \eqref{inf:formal-update} are additions,
multiplications, coefficient extractions and divisions by the
nonzero quantities already listed. The resulting coefficients
are analytic in a common neighbourhood of the compact real
parameter set for each finite jet. The neighbourhood may shrink
with its order. At a fixed order only finitely many collar
normalizations and monic connections occur, so a positive common
neighbourhood exists for each finite order.

Through order $\eps^4$ the seed equation is the expansion of
$\chi a_2+g_2a_2^2/2$. All cubic seed projections start at order
$\eps^5$ by the degree loss just described. Hence the chosen root
agrees with $-2\chi/g_2$ through order $\eps^2$.
The order-$\theta^2\eps^2$ shape coefficient vanishes by the
following cancellation.
At that order the index-four amplitude is $-a_2^2/2$, whose
physical normal derivative tends to $-1$.
The quadratic Neumann correction is
$-(2p-1)a_2^2\mathsf H_2^2/(2R)$ and $(2p-1)/R\to1$.
Their index-four contributions cancel.
All lower indices of $\mathsf H_2^2$ gain at least one $\eps$, and the other
resonant amplitudes start at order $\theta^2\eps^3$.
It follows that $h-a_2\mathsf H_2=O(\theta^2\eps^3)$ at these orders,
proving \eqref{inf:centre-profile} for any sufficiently long finite
jet with an $O_N(|\theta|\eps^3)$ remainder.
A finite truncation of the linear-in-$\theta$ part is included in
this latter bound. Inserting \eqref{inf:chi-analytic} and
\eqref{inf:g2-exact} gives \eqref{inf:centre-witness}.

Choose the recursion order greater than $N+4$ and retain its finite
set of modes. Evaluate the truncated coefficients in the exact
finite-$p$ polynomials $\mathsf H_j$ and exact normalized radial functions.
Restoring the factors $\theta^2\eps^3$ and $\theta^2\eps^4$
in \eqref{inf:seed-equation} and \eqref{inf:other-equations}
recovers the original resonant Dirichlet coefficients. The two
unit eliminations preserve the chosen Taylor jet of the Neumann
and nonresonant Dirichlet equations. Thus a jet of order greater
than $N+4$ makes every original defect coefficient through order
$N$ vanish.
Choose the truncations so that $h-a_2\mathsf H_2$ has a factor
$\theta^2$, using the identical truncated $a_2$ in the field and
the boundary. The higher powers of $\theta$ already present in
$a_2$ are retained in this convention.
For $x$ in any fixed bounded complex neighbourhood $U'\supset U$,
this polynomial has size $O_{N,U'}(|\theta|\eps)$.
For sufficiently small $|\eps|$ its values lie in a fixed collar,
where the integral equation for \eqref{inf:collar-ode} gives jointly
analytic solutions with uniform bounds. The denominators
$\rho+\eps y$ remain away from zero, and only finitely many mode
normalizations occur. Thus both boundary defects are analytic on a
common neighbourhood of the compact parameter set and $U'$.
Their Taylor coefficients in $\eps$ vanish through order $N$ by the
finite induction, as polynomial identities in $x$.

Both defects have an analytic factor $\theta^2$.
At $\theta=0$ the field and domain are the ball. The
linear Neumann defect is
$-\partial_\theta h+(\partial_\theta a_2)\mathsf H_2=0$ by the identical
seed choice; the linear Dirichlet defect is zero since
$u'_0(R)=0$ and $f_2(R)=\chi$ itself has a factor $\theta$.
All other terms have at least two such factors.
Divide the analytic defects by $\theta^2$ and apply Taylor's formula
with remainder in $\eps$ on a smaller fixed complex disk. Compactness
in $r,\theta,x$ bounds the remainder by $C_{N,U}|\eps|^N$.
Cauchy's formula on $U'$ gives every stipulated fixed angular
derivative bound. This proves \eqref{inf:centre-complex-defect} and
the remainder statements. At actual integer blocks the finite regular field is
analytic across axes and origin by Lemma~\ref{inf:solid-harmonics}.
\end{proof}

\begin{lemma}[Analytic defects give weighted coefficients]
\label{inf:analytic-coefficients}
Let $\mu$ be any positive measure on $[0,1]$ with infinite support,
$\mathsf H_j$ its monic orthogonal polynomials, and suppose $0<\mathsf H_j(1)\le1$.
For a complex-valued function $F$ analytic on a neighbourhood of
$U_d=\{z:\dist(z,[0,1])\le d\}$, its coefficients in
$G_j=\mathsf H_j/\mathsf H_j(1)$ satisfy
\begin{equation}\label{inf:cauchy-coefficients}
 |F_j|\le d^{-j}\sup_{U_d}|F|.
\end{equation}
Consequently, for each fixed $k\ge0$ and $1<\rho'<d$,
\begin{equation}\label{inf:weighted-analytic-coefficients}
 \sum_{j\ge0}(1+j)^k(\rho')^j|F_j|
 \le\left(\sum_{j\ge0}(1+j)^k(\rho'/d)^j\right)\sup_{U_d}|F|.
\end{equation}
The constants do not depend on the measure. These assertions apply
to all the beta measures used here.
\end{lemma}
\begin{proof}
For $j=0$ the coefficient is the normalized integral of $F$, so
\eqref{inf:cauchy-coefficients} follows directly. Let $j\ge1$.
The zeros $x_1,\ldots,x_j$ of a monic orthogonal polynomial are simple
and lie in $(0,1)$. Interpolate $F$ at those nodes by a polynomial
$L$ of degree at most $j-1$. The divided-difference identity gives
\[
 F(x)-L(x)=F[x_1,\ldots,x_j,x]\mathsf H_j(x).
\]
The integral formula for divided differences, valid also for
complex-valued functions, gives
\[
 |F[x_1,\ldots,x_j,x]|\le\frac1{j!}
                                  \sup_{[0,1]}|F^{(j)}|.
\]
Orthogonality to $L$ therefore gives
\[
 \frac{|\langle F,\mathsf H_j\rangle|}{\nrm{\mathsf H_j}_{L^2(\mu)}^2}
 \le\frac{\sup_{[0,1]}|F^{(j)}|}{j!}
 \le d^{-j}\sup_{U_d}|F|,
\]
where the last inequality is Cauchy's estimate on each radius-$d$
disk centred on the interval. Multiplication by $\mathsf H_j(1)\le1$
proves \eqref{inf:cauchy-coefficients}; summation gives
\eqref{inf:weighted-analytic-coefficients}.
For our measure \eqref{inf:angular-jacobi} shows
$\mathsf H_j(1)=(\bb)_j/(p+j-1)_j\le1$.
\end{proof}

\begin{corollary}[The coefficient form of the centre defects]
\label{inf:centre-coefficient-defect}
Fix $N,k$ and $\rho'>1$. The centre in Theorem~\ref{inf:centres}
can be chosen so that its physical Dirichlet and radial Neumann
boundary defects $D_N,N_N^{\mathrm{phys}}$, expanded in $G_j$, obey
\begin{equation}\label{inf:centre-coeff-defect-bound}
 \sum_j(1+j)^k(\rho')^j
       (|D_{N,j}|+|N_{N,j}^{\mathrm{phys}}|)
 \le C_{N,k,\rho'}\theta^2p^{-N}.
\end{equation}
The constants and threshold are uniform in $r$ and $0<|\theta|\le20$.
\end{corollary}
\begin{proof}
Choose $d>\rho'$ and then one fixed $R_{\mathrm{an}}>1+d$.
The set $U_d$ lies in the disc of radius $R_{\mathrm{an}}$.
Use the analytic remainder established by the finite-jet construction
in Theorem~\ref{inf:centres} on that disc, and apply
Lemma~\ref{inf:analytic-coefficients} separately to the two defects.
The radius choice changes only the constants and threshold for
this fixed jet. For the Newton application choose $\rho'>2$ once.
\end{proof}

\begin{proposition}[Transfer of an invariant spectral gap]
\label{inf:centre-gap-transfer}
Fix the centre order $N$ and suppose the actual quartic domain has
invariant Dirichlet gap at least $c_{\mathrm q}\tau p^{-5/3}$ about
$R^2$ in unit-ball normalization, with fixed $c_{\mathrm q}>0$.
Then, for all sufficiently large $p$ depending on $N,c_{\mathrm q}$,
the centre domain $\Omega_N/R$ has gap at least
$(c_{\mathrm q}/2)\tau p^{-5/3}$. The threshold is independent of
the admissible split and of $0<\tau\le20$.
\end{proposition}
\begin{proof}
By \eqref{inf:centre-profile}, the two normalized radius functions
differ by $O_N(\tau p^{-4})$ and are both $1+O_N(\tau p^{-2})$.
Thus, for $\delta\le C_N\tau p^{-4}$,
\[
 (1-\delta)\Omega_{\mathrm q}/R\subset\Omega_N/R
 \subset(1+\delta)\Omega_{\mathrm q}/R.
\]
Zero extension preserves the invariant sector, so the min--max
principle gives, for every ordered index $i$,
\[
 (1+\delta)^{-2}\lambda_i(\Omega_{\mathrm q}/R)
 \le\lambda_i(\Omega_N/R)
 \le(1-\delta)^{-2}\lambda_i(\Omega_{\mathrm q}/R).
\]
For any eigenvalue within a fixed window of $R^2$, its counterpart
is $O(p^2)$, and the difference is bounded by $C_N\tau p^{-2}$.
This is smaller than $(c_{\mathrm q}/2)\tau p^{-5/3}$ once
$p^{1/3}>2C_N/c_{\mathrm q}$. An eigenvalue inside the proposed
centre gap would then put its ordered quartic counterpart strictly
inside the original gap, a contradiction. The factor $\tau$ cancels
from the required threshold.
\end{proof}
